\documentclass[11pt,a4paper]{article}
\usepackage[T1]{fontenc}
\usepackage[utf8]{inputenc}
\usepackage{lmodern}
\usepackage{amsmath,amssymb}
\usepackage[margin=25mm]{geometry}
\usepackage{microtype}
\usepackage{hyperref}
\hypersetup{colorlinks=true,linkcolor=blue,citecolor=blue,urlcolor=blue,
 pdftitle={Eventual infinite log-concavity of positive polynomial powers},
 pdfauthor={Aurantium},
 pdfsubject={Fixed positive full-support polynomials, sufficiently large integer powers and all iterations}}
\setlength{\emergencystretch}{3em}
\allowdisplaybreaks
\title{Eventual infinite log-concavity\ of positive polynomial powers}
\author{Aurantium}
\date{14 October 2026}
\begin{document}
\maketitle
\begin{abstract}
For every fixed polynomial whose coefficients are strictly positive throughout its degree range, we prove that the coefficient sequence of every sufficiently large integer power is strictly infinitely log-concave on its original support. The threshold may depend on the polynomial but is independent of the subsequent iteration count. The proof combines uniform exponentially tilted Gaussian interpolation in the bulk with strict Poisson certificates and continuity at the endpoints. We use the established local stopping and shrinking-disk framework, deriving the estimates needed here with explicit conservative constants. For trinomial powers we also obtain an explicit cofinite exponent threshold and a uniform logarithmic-curvature estimate through a growing number of iterations. The sharper assertion that every trinomial exponent at least ten has the property remains unresolved.
\end{abstract}

\section{Introduction and main results}

For a zero-extended sequence \(a=(a_k)\), the log-concavity operator is
\[
 (\mathcal L a)_k=a_k^2-a_{k-1}a_{k+1}.
\]
A sequence is infinitely log-concave if all its iterates under \(\mathcal L\) are entrywise nonnegative. We say that the iterates are strictly positive on the original support when every coefficient in that support is positive at every stage; the prescribed exterior zeros remain zero.

Medina and Straub formulate a sufficiently-large-power conjecture for polynomials with positive coefficients throughout their degree range \cite[Conjecture~5.2]{MS}. The result here addresses its integer-exponent formulation. The polynomial is fixed, its coefficients may be arbitrary positive real numbers, and no threshold uniform over all polynomials is asserted.

\textbf{Theorem B (eventual positive-polynomial convolution theorem).}
Let \[
 p(x)=\sum_{i=0}^d c_i x^i,\qquad c_i>0\quad(0\leq i\leq d).
\] There exists \(N(p)\) such that, for every integer \(n\geq N(p)\),
the coefficient sequence of \(p(x)^n\), with zero extension, is strictly
infinitely log-concave on its positive support.

This is the main result of the paper. Its threshold \(N(p)\) works simultaneously for every subsequent iteration; it is not a threshold chosen separately for each iteration. Degree zero is included as the positive singleton case. No claim about noninteger powers, internal zero coefficients or arbitrary coefficient signs is made.

The analytic tools are the local stopping and shrinking-disk framework developed by Xie, Zhang and Zhang \cite[Section~2]{XZZ}, which also credits related earlier work of Liu and Mao. Sections~2--3 give complete proofs of the forms used here, with conservative explicit constants. The application to polynomial powers requires a tilted interpolation uniform near the endpoints and a transfer of finite strict Poisson certificates. General convolution closure for infinite log-concavity is not assumed; Medina--Straub's product example shows why that shortcut is unavailable \cite[Example~5.1]{MS}. The factor-cone and determinant background for the accompanying finite computations is discussed in \cite{MCS}.

The trinomial polynomial \(1+x+x^2\) supplies an explicit companion result. We retain the following notation and distinguish the eventual conclusion from the sharper threshold problem.

For integers \(n\geq0\) and \(k\in\mathbb Z\), put \[
 T(n,k)=[x^k](1+x+x^2)^n,
 \qquad T(n,k)=0\quad(k<0\text{ or }k>2n),
\] and define, retaining this zero extension, \[
 a^{(0)}_{n,k}=T(n,k),\qquad
 a^{(j+1)}_{n,k}=(a^{(j)}_{n,k})^2-a^{(j)}_{n,k-1}a^{(j)}_{n,k+1}.
\] The original target remains \[
 \boxed{\forall n\in\mathbb Z_{\geq10}\;\forall j\in\mathbb Z_{\geq0}\;
        \forall k\in\mathbb Z:\ a^{(j)}_{n,k}\geq0.} \tag{T}
\]
This is the explicit-threshold assertion of Medina--Straub \cite[Example~5.3]{MS}. It retains every integer exponent \(n\geq10\), every iteration and every coefficient, with zero extension and weak inequalities. It remains unresolved in this paper.

\textbf{Theorem A (explicit cofinite-exponent theorem).} Define \[
 \mathcal N=2^{\,2^{40100}}.
\] For every integer \(n\geq\mathcal N\), every iterate is strictly
positive on the original support: \[
 a^{(j)}_{n,k}>0\qquad(j\geq0,\ 0\leq k\leq2n).
\] Moreover, at the common stage \[
 J(n)=\left\lceil\log_2(8n)\right\rceil
\] the entire row is 3-factor log-concave, and this remains true at
every later stage. A smaller explicit cutoff \(N_0\), defined in Section
7, also works.

The bound \(\mathcal N\) is derived symbolically from analytic and
polynomial inequalities. \textbf{It was neither evaluated as an integer
nor verified by enumerating exponents up to it.} The proof controls
growing iteration depth, not a fixed number of iterations followed by an
unjustified limit.

Theorem~C, stated and proved in Section~9, gives uniform normalized curvature through the growing horizon \(\lceil\log_2(6V)\rceil\) for trinomial rows, where \(V\) is the tilted variance of Section~5. At the center this horizon is proportional to \(\log n\). The result also gives local cone-entry bounds; its positivity statements concern the required local dependency intervals. The proof of the main Theorem~B is in Section~8, following the analytic and Poisson ingredients in Sections~2--7.

\section{Ratios, curvature, and a local-to-global stopping argument}\label{ratios-curvature-and-a-local-to-global-stopping-argument}


For a positive sequence on a finite interval or on
\(\mathbb Z_{\geq0}\), define its interior ratios and logarithmic
curvatures by \[
 q_j(k)=\frac{a^{(j)}_{k-1}a^{(j)}_{k+1}}{(a^{(j)}_k)^2},
 \qquad h_j(k)=-\log q_j(k).
\] At an original zero-extension endpoint set \(q_j=0\). Curvatures are
used only at interior indices with positive neighbors. If all needed
ratios lie in \([0,1)\), direct substitution gives \[
 q_{j+1}(k)=
 \frac{q_j(k)^2(1-q_j(k-1))(1-q_j(k+1))}{(1-q_j(k))^2}. \tag{1}
\] Let \(\Delta f(z)=f(z-1)-2f(z)+f(z+1)\). Then the curvature
recurrence is \[
 h_{j+1}=2h_j-\Delta g(h_j),\qquad
 g(w)=\operatorname{Log}(1-e^{-w}). \tag{2}
\] On \(\operatorname{Re}w>0\), \(g\) is holomorphic and
\(|\operatorname{Im}g(w)|<\pi/2\): indeed,
\(\operatorname{Re}(1-e^{-w})>0\).

\textbf{Lemma 1 (local stopping implies global positivity).} Suppose
that for each interior index \(k\) there is a finite \(T(k)\) such that
its actual coefficient iterates are positive through \(T(k)\) and its
ratio at that time is defined using positive neighbors and satisfies \[
 q_{T(k)}(k)\leq\frac13.
\] The local statements may be established on separate finite dependency
intervals. Then all coefficient iterates are positive everywhere on the
original support. Each attained ratio bound persists at every later
stage. On finite support, the entire row satisfies the bound at stage
\(\max_kT(k)\).

\textbf{Proof.} When the ratios needed in (1) lie in \([0,1)\), \[
 q_j(k)\leq\tfrac13
 \quad\Longrightarrow\quad
 q_{j+1}(k)\leq\frac{q_j(k)^2}{(1-q_j(k))^2}
 \leq\tfrac14<\tfrac13. \tag{3}
\] Suppose a first nonpositive layer \(r\) exists and contains a
nonpositive entry at \(k\). Endpoints are repeatedly squared, so \(k\)
is interior. Its local certificate gives \(T(k)<r\). All entire layers
before \(r\) are positive, so all ratios needed to propagate (3) through
layer \(r-1\) lie in \([0,1)\). Thus \(q_{r-1}(k)\leq1/3\), and \[
 a^{(r)}_k=(a^{(r-1)}_k)^2(1-q_{r-1}(k))>0,
\] a contradiction. This least-layer argument also works on the
half-line; it does not require taking a maximum over infinitely many
stopping times. Persistence then follows from (3). \(\square\)

This is a specialization of the local criterion in Xie--Zhang--Zhang
{[}2, Lemma 2.1{]}. It is important that no convolution-closure property
for infinite log-concavity is assumed.

Positive global and geometric rescaling preserves all ratios.
Explicitly, if \(b_k=AB^ka_k\) with \(A,B>0\), then \[
 (\mathcal L^j b)_k=A^{2^j}B^{k2^j}(\mathcal L^j a)_k. \tag{4}
\]

\section{Holomorphic control through growing iteration depth}\label{holomorphic-control-through-growing-iteration-depth}


The next lemma is a conservative, self-contained version of the
shrinking-disk estimates in {[}2, Section 2{]}. Its constants are chosen
for transparent arithmetic rather than sharpness.

\textbf{Lemma 2 (holomorphic amplification).} Let \(h_0\) be holomorphic
on a neighborhood of the closed disk \(|z-k|\leq R\), real on its real
diameter, and suppose \[
 \operatorname{Re}h_0\geq\mu,\qquad
 0<\mu\leq1,\qquad \mu R^2\geq256. \tag{5}
\] Put \[
 K=\left\lceil\log_2(4/\mu)\right\rceil,
 \qquad R_j=R(3/4)^j\quad(0\leq j\leq K).
\] The recurrence (2) defines holomorphic functions with positive real
part on the corresponding disks. For \(1\leq j\leq K\), \[
 \left|2^{-j}h_j(z)-h_0(z)\right|<\frac{56}{R^2}
 \quad (|z-k|\leq R_j), \tag{6}
\] and \[
 h_K(k)>\frac{25}{8}>\log3. \tag{7}
\] If \(h_0\) agrees with the curvature of positive initial coefficients
at all required integer samples, these functions agree with the
curvatures of positive local coefficient iterates; thus \(q_K(k)<1/3\).

\textbf{Proof.} First let \(F\) be holomorphic on a radius-\(d\) disk
about \(z\), with \(|\operatorname{Im}F|\leq\pi/2\). If
\(F(z+w)=\sum_{m\geq0}c_mw^m\), Fourier coefficients on smaller circles
give, for \(m\geq1\), \[
 c_m\rho^m=\frac{i}{\pi}\int_0^{2\pi}
 \operatorname{Im}F(z+\rho e^{it})e^{-imt}\,dt.
\] After rotating to extract the absolute value, the identity
\(\int_0^{2\pi}|\cos(mt+\alpha)|\,dt=4\) shows \(|c_m|\leq2/d^m\) by
passage to the boundary. Consequently, for \(d>1\), \[
 |\Delta F(z)|\leq2\sum_{m\geq1}|c_{2m}|
 \leq\frac4{d^2-1}. \tag{8}
\]

Whenever the earlier curvatures have positive real part, telescoping (2)
and applying (8) on the nested disks yields \[
 |2^{-j}h_j-h_0|\leq E_j,
 \qquad
 E_j=\sum_{i=0}^{j-1}
 \frac{2^{1-i}}{(R_i-R_j)^2-1}. \tag{9}
\] We verify that all denominators are positive and \(E_j<56/R^2\)
without assuming the induction conclusion.

Write \(\beta=\log_2(4/3)<1/2\). The smallest successive radius gap is
\(d_{\min}=(R/4)(3/4)^{K-1}\). Since \(K-1<\log_2(4/\mu)\), (5) gives \[
 d_{\min}^2>
 \frac{R^2}{16}(\mu/4)^{2\beta}
 \geq16\mu^{2\beta-1}4^{-2\beta}
 \geq\frac{81}{16}>4. \tag{10}
\] For \(b=8/9\), it follows that \[
 E_j<\frac{8}{3R^2}
 \sum_{r=1}^j\frac{b^{j-r}}{(1-(3/4)^r)^2}.
\] The reciprocal denominators are at most \(16,6,3,7/3\) for
\(r=1,2,3,r\geq4\), respectively. The sum is less than 21: check the
first three cases using \(16b^2+6b+3<21\), and then use its recurrence
and \((7/3)/(1-b)=21\). Therefore \(E_j<56/R^2\).

Induction in (9) now gives \[
 \operatorname{Re}h_j>
 2^j(\mu-56/R^2)\geq\frac{25}{32}2^j\mu>0.
\] Since \(2^K\mu\geq4\), (7) follows. Holomorphic logarithms are
legitimate at every step because these real parts are positive.

For the sequence interpretation, every radius gap exceeds two. The real
integer samples and their unit shifts needed for a new curvature
therefore lie in the previous domain. Positive initial entries and
positive real curvatures give positive next coefficients by
\(a^{(j+1)}=(a^{(j)})^2(1-e^{-h_j})\). Induction identifies the
recurrence with the actual coefficient iterates on successively smaller
index intervals, including positive neighboring entries at the final
center. No artificial boundary values are assigned. \(\square\)

The key feature is the uniform normalized error in (6), despite \(K\)
increasing as \(\mu\) decreases. Merely applying a fixed-depth normal
approximation would not give this conclusion.

\section{A completely certified Poisson boundary model}\label{a-completely-certified-poisson-boundary-model}


Let \(u_k=1/k!\) for \(k\geq0\), and \(u_k=0\) for \(k<0\).

\textbf{Lemma 3 (strict local cone entry for the Poisson seed).} For
each \(k\geq1\) there is a local certificate with a stopping time \[
 T_P(k)\leq\left\lceil\log_2(8k)\right\rceil
\] such that the required iterates are positive through that time and
\(q_{T_P(k)}(k)<1/3\). Consequently all iterates of \(u\) are strictly
positive on the half-line.

\textbf{Proof for \(k\geq2048\).} Its initial curvature has the
interpolation \[
 h_0(z)=\operatorname{Log}(1+1/z)=\int_0^1\frac{dt}{z+t}.
\] On \(|z-k|\leq k/2\), inversion of the real-centered disk gives \[
 \operatorname{Re}\frac1{z+t}\geq\frac1{3k/2+t}.
\] For completeness, if \(w\) belongs to \(|w-a|\leq b<a\), the image
under \(w\mapsto1/w\) is a disk whose leftmost real point is
\(1/(a+b)\); this proves the displayed inequality with \(a=k+t\),
\(b=k/2\). Hence \[
 \operatorname{Re}h_0(z)\geq\frac1{3k/2+1}\geq\frac1{2k}.
\] Apply Lemma 2 with \(\mu=1/(2k)\) and \(R=k/2\), since
\(\mu R^2=k/8\geq256\). The disks and all needed initial unit shifts lie
in the positive half-line. This gives the stated time bound and a strict
terminal gap.

\textbf{Exact finite certificate for \(1\leq k\leq2047\).} Initially
\(q_0(k)=k/(k+1)\), including \(q_0(0)=0\). The verifier uses
denominator \(D=2^{256}\) and outward-rounded dyadic intervals. If
\(l_k\leq Dq_k\leq v_k<D\), the next bounds are \[
\begin{split}
 l'_k&=\left\lfloor
 \frac{l_k^2(D-v_{k-1})(D-v_{k+1})}{D(D-l_k)^2}
 \right\rfloor,\\
 v'_k&=\left\lceil
 \frac{v_k^2(D-l_{k-1})(D-l_{k+1})}{D(D-v_k)^2}
 \right\rceil.
\end{split} \tag{11}
\] They are valid because (1) increases with its central argument and
decreases with either neighbor on \([0,1)\). All denominator and
positivity inequalities are checked with integers.

The initial ratio prefix ends at index 2064. The last available index is
discarded at each step, never declared to have a zero neighbor. Every
index \(1\leq k\leq2047\) obtains a strict certificate \(3v_k<D\) by
depth 12. The numbers certified by stages \(0,\ldots,12\) are \[
 0,0,2,6,13,27,56,115,231,464,930,1862,2047.
\] Each first certified time also satisfies
\(T_P(k)\leq\lceil\log_2(8k)\rceil\). Every strict gap and its exact
dyadic bounds are in \texttt{data/poisson\_certificates.json.gz}.

An independent check starts with the integer prefix \(64!/k!\), applies
the coefficient definition of \(\mathcal L\) directly, again discarding
the artificial right edge, and compares its exact ratios with the
interval bounds through eight iterations. All 540 comparisons pass. The
computed finite certificates and the analytic certificates for
\(k\geq2048\), followed by Lemma 1, finish the proof. \(\square\)

The finite computations here are proof certificates, not extrapolation
from experiments. They are reconstructed from the displayed rational
formulas when the verifier is run.

\section{Uniform exponentially tilted Gaussian approximation}\label{uniform-exponentially-tilted-gaussian-approximation}


For \(1\leq k\leq n\), let \(\rho\in(0,1]\) be the unique solution of \[
 m(\rho)=\frac{\rho+2\rho^2}{1+\rho+\rho^2}=\frac{k}{n}.
\] Set \(p(\rho)=1+\rho+\rho^2\) and \[
 v(\rho)=\frac{\rho(1+4\rho+\rho^2)}{p(\rho)^2},\qquad V=nv(\rho).
\] Existence and uniqueness follow because \(m\) increases continuously
from 0 to 1 on \((0,1]\) and \(\rho m'(\rho)=v(\rho)>0\).

The elementary variance bounds used below are \[
 \frac23k\leq V\leq\frac98k. \tag{12}
\] Indeed, \(v/m=(1+4\rho+\rho^2)/(p(\rho)(1+2\rho))\). For the lower
inequality the numerator difference is \[
 3(1+4\rho+\rho^2)-2p(\rho)(1+2\rho)
 =(1-\rho)(4\rho^2+7\rho+1)\geq0.
\] For the upper inequality it is \[
 9p(\rho)(1+2\rho)-8(1+4\rho+\rho^2)
 =1-5\rho+19\rho^2+18\rho^3>0,
\] since \(1-5\rho+19\rho^2=19(\rho-5/38)^2+51/76>0\).

Let \(X_\rho\) have probabilities \((1,\rho,\rho^2)/p(\rho)\) at
\(0,1,2\). Its mean and variance are \(m(\rho)\) and \(v(\rho)\). Put \[
 \phi_\rho(t)=\mathbb E e^{it(X_\rho-m(\rho))}
 =e^{-im(\rho)t}\frac{p(\rho e^{it})}{p(\rho)}
\] and introduce the entire function \[
 F_{n,k}(z)=\frac{\sqrt V}{2\pi}\int_{-\pi}^{\pi}
 \phi_\rho(t)^n e^{-i\sqrt Vzt}\,dt. \tag{13}
\] At every integer \(s\), Fourier orthogonality gives the exact
sampling identity \[
 F_{n,k}(s/\sqrt V)=
 \frac{\sqrt V\,\rho^{k+s}}{p(\rho)^n}\,T(n,k+s). \tag{14}
\] The factors on the right are positive global and geometric factors;
by (4) they do not change iterated ratios. The tilt \(\rho\) and
variance \(V\) are held fixed across this entire local interval and all
its subsequent local iterations.

\textbf{Lemma 4 (explicit uniform analytic normal bound).} For every
\(A\geq0\), \[
 \sup_{|z|\leq A}
 \left|F_{n,k}(z)-G(z)\right|
 \leq\frac{64e^{8A^2}}{\sqrt V},
 \qquad G(z)=\frac1{\sqrt{2\pi}}e^{-z^2/2}. \tag{15}
\] The bound is uniform in all the allowed \(n,k,\rho\), including
\(\rho\to0\).

\textbf{Proof.} Write \(\pi_i=\mathbb P(X_\rho=i)\). The contribution of
adjacent pairs gives \[
 |\phi_\rho(t)|^2\leq
 1-2(\pi_0\pi_1+\pi_1\pi_2)(1-\cos t).
\] Since \[
 \pi_0\pi_1+\pi_1\pi_2
 =\frac{\rho(1+\rho^2)}{p(\rho)^2}\geq\frac{v(\rho)}3,
\] we get, for \(|t|\leq\pi\), \[
 |\phi_\rho(t)|\leq
 \exp\{-v(\rho)(1-\cos t)/3\}
 \leq e^{-v(\rho)t^2/16}. \tag{16}
\] The last inequality uses \(1-\cos t\geq2t^2/\pi^2\) and \(\pi^2<10\).

For \(Y=X_\rho-m(\rho)\), \(|Y|\leq2\), \(\mathbb E|Y|^3\leq2v\), and
\(v\leq1\). Taylor's formula on the real axis gives \[
 \phi_\rho(t)=1-vt^2/2+R(t),\qquad |R(t)|\leq v|t|^3/3.
\] For \(|t|\leq1/4\), let \(w=\phi_\rho(t)-1\). Then
\(|w|\leq(7/12)vt^2\leq7/192\). The logarithm series remainder satisfies
\[
 |\operatorname{Log}(1+w)-w|
 \leq\frac{|w|^2}{2(1-|w|)}
 \leq\tfrac15v^2t^4.
\] Consequently, \[
 |\operatorname{Log}\phi_\rho(t)+vt^2/2|\leq v|t|^3. \tag{17}
\]

Change variables to \(u=\sqrt Vt\) in (13), extending its integrand by
zero outside \([-\pi\sqrt V,\pi\sqrt V]\). When \(|u|\leq\sqrt V/4\),
(17) gives an error at most \(|u|^3/\sqrt V\) in the exponent, and both
real parts in the exponential comparison are at most \(-u^2/4\). Thus \[
 |\phi_\rho(u/\sqrt V)^n-e^{-u^2/2}|
 \leq\frac{|u|^3}{\sqrt V}e^{-u^2/4}. \tag{18}
\] When \(|u|>\sqrt V/4\), (16) and the Gaussian bound instead give \[
 |\phi_\rho(u/\sqrt V)^n\mathbf1_{|u|\leq\pi\sqrt V}-e^{-u^2/2}|
 \leq\frac{8|u|}{\sqrt V}e^{-u^2/16}. \tag{19}
\] For complex \(|z|\leq A\), multiply these bounds by \(e^{A|u|}\).
Young's inequality and direct integration yield \[
 \int_{\mathbb R}|u|^3e^{-u^2/4+A|u|}\,du\leq64e^{2A^2},
\] \[
 \int_{\mathbb R}8|u|e^{-u^2/16+A|u|}\,du\leq256e^{8A^2}.
\] The Fourier integral of \(e^{-u^2/2}\) is \(\sqrt{2\pi}e^{-z^2/2}\);
the identity extends to complex \(z\) by holomorphy. Dividing the last
two bounds by \(2\pi\sqrt V\) proves (15), since \(320/(2\pi)<64\).
\(\square\)

\section{Uniform bulk certificates for all iterations}\label{uniform-bulk-certificates-for-all-iterations}


\textbf{Proposition 5 (a local certificate at every large-variance
position).} If \(V\geq2^{20000}\), then the actual trinomial iterates
near \(k\) have a strict local 3-factor certificate at stage \[
 T(n,k)=\left\lceil\log_2(6V)\right\rceil
 \leq\left\lceil\log_2(8k)\right\rceil. \tag{20}
\] The positivity statement includes the entire finite dependency
interval needed to define this local ratio. It requires no assumption
about positivity of distant entries at these stages.

\textbf{Proof.} Apply Lemma 4 on \(|z|\leq24\). Since
\(|G(z)|\geq(2\pi)^{-1/2}e^{-288}\) there, \[
 |F_{n,k}/G-1|
 \leq\frac{192e^{4896}}{\sqrt V}<2^{-200}<\tfrac14. \tag{21}
\] For the strict inequality, use \(192<2^8\), \(e<4\), and
\(\sqrt V\geq2^{10000}\); then \(10000-2\cdot4896-8=200\).

The function \(\eta=\operatorname{Log}(F_{n,k}/G)\) is holomorphic near
that disk and satisfies \(|\eta|<1/2\). It is real on the real diameter,
because the integral defining \(F_{n,k}\) is real there and (21) makes
\(F_{n,k}/G\) positive. By Cauchy's estimate on radius-three disks, \[
 |\eta''(z)|\leq1/9\qquad(|z|\leq21).
\] Put \(\delta=1/\sqrt V\) and, for \(w=k+\sqrt Vz\), define the
physical-index curvature interpolation \[
 H_0(w)=2\log F_{n,k}(z)-\log F_{n,k}(z-\delta)
                       -\log F_{n,k}(z+\delta).
\] Use the logarithm \(\log G+\eta\) here. The Gaussian contributes
exactly \(\delta^2\), and the integral formula for a symmetric second
difference gives \[
 |VH_0(w)-1|\leq\frac19
 \quad\text{when }|w-k|\leq20\sqrt V. \tag{22}
\] Thus (5) holds with \[
 R=20\sqrt V,\qquad \mu=\frac2{3V},\qquad
 \mu R^2=\frac{800}{3}>256.
\] Lemma 2 proves the certificate at \(K=\lceil\log_2(6V)\rceil\).

There is no support problem: (12) implies \(k\geq V/2\), and
\(V\geq2^{20000}>6400\) implies \(20\sqrt V\leq k/2\). The initial disk
and its required unit shifts lie strictly inside \([0,2n]\) for these
parameters. Equation (14) identifies its integer values with positive
scaled trinomial coefficients, so Lemma 2 applies to the actual row.
Finally \(6V\leq27k/4<8k\), proving the time bound. \(\square\)

By symmetry, the same conclusion holds near \(2n-k\). In particular, set
\[
 B=2^{20001}.
\] Every position with \[
 \min(k,2n-k)\geq B \tag{23}
\] has a certificate as in Proposition 5, uniformly in \(n\), because
\(V\geq(2/3)B>2^{20000}\). The remaining problem is a fixed-width
boundary, not an uncontrolled region growing proportionally to \(n\).

\section{Effective boundary control and proof of Theorem A}\label{effective-boundary-control-and-proof-of-theorem-a}


For \(m\geq0\) define the rational polynomial \[
 P_{0,m}(x)=\sum_{j=0}^{\lfloor m/2\rfloor}
 \frac{(x)_{\underline{m-j}}}{j!(m-2j)!}. \tag{24}
\] For every nonnegative integer \(n\), \(P_{0,m}(n)=T(n,m)\), including
\(m>2n\), when the polynomial evaluates to zero. This follows by
counting the terms with \(j\) choices of \(x^2\) and \(m-2j\) choices of
\(x\).

Set \(P_{\ell,-1}=0\) and recursively define \[
 P_{\ell+1,m}=P_{\ell,m}^2-P_{\ell,m-1}P_{\ell,m+1}. \tag{25}
\] These are exactly the trinomial iterate polynomials evaluated at
\(n\).

\textbf{Lemma 6 (leading coefficients).} If
\(u^{(\ell)}=\mathcal L^\ell(1/m!)_{m\geq0}\), then \[
 P_{\ell,m}(x)=u^{(\ell)}_m x^{m2^\ell}+\text{lower-degree terms}. \tag{26}
\] In particular, the leading coefficient is strictly positive.

\textbf{Proof.} The leading coefficient of (24) is \(1/m!\). Both terms
in (25) have degree at most \(m2^{\ell+1}\), and their
leading-coefficient difference is exactly \(u^{(\ell+1)}_m\). Lemma 3
shows this number is positive, so it cannot cancel. Induction proves
(26). \(\square\)

We make all dependence on the finite boundary explicit. Let \[
 J=20004,\qquad M=B+J+1=B+20005,\qquad
 S=M!,\qquad H=2^{M+1}S,
\] and define \[
 \boxed{N_0=2H^{\,2^{20005}}.} \tag{27}
\] Only polynomial indices \(m\leq M-\ell\) will be used at layer
\(\ell\). This convention discards insufficiently supported right-edge
calculations rather than treating their missing neighbors as zero.

Write \(\|f\|_1\) for the sum of the absolute values of the coefficients
of a polynomial. In (24), the norm of \((x)_{\underline t}\) is \(t!\)
(including \(t=0\)). Hence \[
 \|P_{0,m}\|_1\leq
 \sum_{j=0}^{\lfloor m/2\rfloor}\binom{m-j}{j}\leq2^m\leq2^M. \tag{28}
\] The middle sum is bounded by \(2^m\) directly from the recurrence of
these binomial sums, or by counting compositions of \(m\) with parts 1
and 2. Since \(\|fg\|_1\leq\|f\|_1\|g\|_1\), (25) gives \[
 \|P_{\ell,m}\|_1\leq
 2^{(M+1)2^\ell-1}. \tag{29}
\] Also \(S^{2^\ell}P_{\ell,m}\in\mathbb Z[x]\): each denominator in
(24) divides \(M!\), and induction in (25) proves the claim. Thus every
positive leading coefficient in (26) is at least \(S^{-2^\ell}\).

For each \(1\leq k<B\), choose the Poisson certificate time
\(T=T_P(k)\leq J\) from Lemma 3 and put \[
 G_k(x)=P_{T,k}(x)^2-3P_{T,k-1}(x)P_{T,k+1}(x). \tag{30}
\] Its leading coefficient is positive, since it is exactly
\((u_k^{(T)})^2-3u_{k-1}^{(T)}u_{k+1}^{(T)}>0\). Its denominator divides
\(S^{2^{T+1}}\), and \[
 \|G_k\|_1\leq4\left(2^{(M+1)2^T-1}\right)^2
 =2^{(M+1)2^{T+1}}. \tag{31}
\] All polynomials needed to define (30), and all its local coefficient
iterates through time \(T\), are among the retained indices:
\(k+T+1\leq M\).

For every required \(P_{\ell,m}\) or \(G_k\), (29)--(31) show \[
 \frac{\|f\|_1}{\text{positive leading coefficient of }f}
 \leq H^{2^{J+1}}. \tag{32}
\] An elementary polynomial sign bound completes the boundary proof. If
\(f(x)=cx^d+\sum_{i<d}c_ix^i\), \(c>0\), and \(x\geq1\), then \[
 f(x)\geq x^{d-1}\left(cx-\sum_{i<d}|c_i|\right)>0
 \quad\text{when }x>\|f\|_1/c.
\] Positive constant polynomials need no bound. Therefore every integer
\(n\geq N_0\) makes all the required boundary entries and all strict
gaps (30) positive. Here \(N_0>M\), so every retained boundary index
lies inside the original support. This gives local certificates at every
\(1\leq k<B\) and, by reflection, at the other boundary. Proposition 5
handles the remaining positions; endpoints are fixed at 1. Lemma 1
proves strict positivity for \textbf{all iterations}.

Every chosen boundary time is at most \(\lceil\log_2(8k)\rceil\), and
every bulk time satisfies (20). After reflection all are at most
\(J(n)=\lceil\log_2(8n)\rceil\). Lemma 1 therefore gives the common
3-factor stage asserted in Theorem A.

Finally, the convenient cutoff \(\mathcal N\) is indeed larger than
\(N_0\). Since \(M<2^{20002}\), \[
 \log_2 H=M+1+\log_2(M!)
 \leq M+1+M\log_2M<2^{20018}.
\] Here one may use \(20003<2^{15}\); the slack is more than sufficient.
Thus \[
 \log_2N_0=1+2^{20005}\log_2H<2^{40024}<2^{40100},
\] so \(N_0<\mathcal N\). This completes the proof of Theorem A. No
enormous polynomial or exponent interval has been enumerated.
\(\square\)

\section{Proof of the main theorem for positive polynomial powers}\label{proof-of-the-general-positive-polynomial-consequence}


We give the proof to separate a genuinely complete unspecified-threshold
result from the unresolved sharper threshold in (T).

\textbf{Proof of Theorem B.} Degree zero is the positive singleton case.
Suppose \(d\geq1\) and \(c_i>0\) for every \(i\) in the degree range.
For \(\rho>0\), let \[
 \mathbb P(X_\rho=i)=\pi_i(\rho)
 =\frac{c_i\rho^i}{p(\rho)},\qquad
 m(\rho)=\mathbb E X_\rho,\qquad v(\rho)=\operatorname{Var}X_\rho.
\] Then \(v(\rho)>0\) and \(\rho m'(\rho)=v(\rho)\), so \(m\) maps
\((0,\infty)\) bijectively onto \((0,d)\). For each interior index \(k\)
choose \(nm(\rho)=k\), and put \(V=nv(\rho)\).

The ratio \[
 \frac{v(\rho)}{\min(m(\rho),d-m(\rho))}
\] is continuous and positive, and tends to 1 at both ends \(\rho\to0\)
and \(\rho\to\infty\). The endpoint limits use
\(c_0,c_1,c_{d-1},c_d>0\); for instance, near zero,
\(m(\rho)=(c_1/c_0)\rho+O_p(\rho^2)\) and
\(v(\rho)=(c_1/c_0)\rho+O_p(\rho^2)\). Compactness after adjoining the
endpoints gives constants \(0<c_p\leq C_p<\infty\) with \[
 c_p b\leq V\leq C_p b,
 \qquad b=\min(k,nd-k). \tag{33}
\]

There is also a uniform characteristic-function bound \[
 |\mathbb E e^{it(X_\rho-m(\rho))}|
 \leq e^{-c'_pv(\rho)t^2},\qquad |t|\leq\pi, \tag{34}
\] where \(c'_p>0\). To prove it for \(0<\rho\leq1\), note that
\(\pi_0(\rho)\pi_1(\rho)/v(\rho)\) is positive continuous and extends to
the positive limit 1 at zero. The adjacent pair \((0,1)\) in the
modulus-squared identity therefore gives (34), using
\(1-\cos t\geq2t^2/\pi^2\). For \(\rho\geq1\), apply the same argument
to the reversed polynomial \(x^dp(1/x)\) with tilt \(1/\rho\). Reversal
leaves the variance and the characteristic-function modulus unchanged.

Furthermore \(|X_\rho-m(\rho)|\leq d\), so its third absolute centered
moment is at most \(d v(\rho)\). The Taylor-logarithm argument used for
(17), with a fixed sufficiently small \(t_p>0\), gives \[
 \left|\log\phi_\rho(t)+v(\rho)t^2/2\right|
 \leq C'_pv(\rho)|t|^3\qquad(|t|\leq t_p), \tag{35}
\] uniformly in \(\rho\). The variance is bounded by \(d^2\), which
makes the logarithm neighborhood uniform. In the scaled variable
\(u=\sqrt Vt\), take \(|u|\leq\min(t_p,1/(4C'_p))\sqrt V\); then (35)
gives an integrable comparison by a constant times
\(|u|^3e^{-u^2/4}/\sqrt V\). On the complement, (34) and the Gaussian
bound give a constant times \(|u|e^{-c''_pu^2}/\sqrt V\) for a fixed
\(c''_p>0\). Multiplication by \(e^{A|u|}\) remains integrable for each
fixed \(A\).

Consequently the tilted Fourier interpolation, defined exactly as in
(13) with this \(p\), satisfies \[
 \sup_{|z|\leq A}|F_{n,k}(z)-G(z)|\leq C_{p,A}/\sqrt V \tag{36}
\] for some finite \(C_{p,A}\) independent of \(n,k,\rho\). The same
exact sampling identity as (14) holds. Taking \(A=24\), then a
sufficiently large threshold \(V_p\), reproduces (21)--(22). Lemma 2
supplies local strict 3-factor certificates whenever \(b\geq B_p\) for
some finite \(B_p\) depending only on \(p\). Enlarge \(B_p\) if
necessary so all physical disks of radius \(20\sqrt V\) and their unit
shifts fit inside the support, using (33). Their times are
\(\lceil\log_2(6V)\rceil=O_p(\log b)\).

It remains to handle finitely many fixed indices near either endpoint.
For every fixed \(m\geq0\), multinomial expansion gives \[
 c_0^{-n}\left(\frac{c_0}{nc_1}\right)^m[x^m]p(x)^n
 \longrightarrow\frac1{m!}\quad(n\to\infty). \tag{37}
\] Indeed, after removing \(c_0^n\), the coefficient is a polynomial in
\(n\) of degree \(m\), with leading coefficient \((c_1/c_0)^m/m!\);
every use of a power \(x^i\) with \(i\geq2\) lowers the number of
selected nonconstant factors.

For a fixed boundary index \(k<B_p\), choose its finite strict Poisson
certificate from Lemma 3. Convergence (37) holds simultaneously on its
entire finite initial dependency interval. The operator \(\mathcal L\)
and the expression \(a_k^2-3a_{k-1}a_{k+1}\) are polynomial functions of
this finite set of entries. Strict positivity of all the finitely many
seed inequalities therefore persists for the rescaled coefficients of
\(p^n\) once \(n\geq N(p,k)\). The rescaling in (37) has the form (4),
so it preserves the signs and ratios being certified. The support
contains this fixed interval for sufficiently large \(n\). Apply the
same argument to the reversed polynomial for the right endpoint.

There are finitely many boundary centers. Take a maximum of their
finitely many thresholds and the support thresholds, and call it
\(N(p)\). All remaining centers have the uniform bulk certificates, and
endpoints are positive powers of \(c_0\) and \(c_d\). Lemma 1 now gives
every iteration on the whole support, simultaneously for each integer
\(n\geq N(p)\). This proves Theorem B. \(\square\)

Every threshold in this proof is allowed to depend on the fixed
polynomial \(p\), exactly as in the integer-exponent version of
Conjecture 5.2. None depends on the final iteration count. The argument
is not an application of an unproved convolution-closure assertion.

\section{Uniform curvature asymptotics at growing depth}\label{uniform-curvature-asymptotics-at-growing-depth}


The new analytic estimates also yield a sharper description of the
iteration scale in the bulk. Use reflection when \(k>n\), and let \(V\)
be the tilted variance associated to the reflected index.

\textbf{Theorem C.} As \(b=\min(k,2n-k)\to\infty\), uniformly in the
allowed \(n,k\), \[
 \sup_{0\leq j\leq\lceil\log_2(6V)\rceil}
 \left|V2^{-j}h_j(k)-1\right|\longrightarrow0. \tag{38}
\] The ratios and curvatures here belong to positive local coefficient
iterates through the indicated depths.

\textbf{Proof.} Fix any \(A\geq20\). Apply Lemma 4 on the scaled disk of
radius \(A+4\). Relative error tends to zero uniformly as
\(V\to\infty\), since the fixed Gaussian has no zeros there. The
analytic logarithm and Cauchy estimates used in Section 6 then give \[
 \sup_{|w-k|\leq A\sqrt V}|VH_0(w)-1|\longrightarrow0.
\] For sufficiently large \(V\), apply Lemma 2 with \(R=A\sqrt V\) and
\(\mu=2/(3V)\); its stopping depth remains \(\lceil\log_2(6V)\rceil\).
Equation (6) implies that the limsup of the left side of (38) is at most
\(56/A^2\). Equation (12) ensures the needed disks lie inside the
support for sufficiently large \(V\). Since \(A\) can be arbitrarily
large, the limsup is zero. This order of choices does not interchange a
fixed-iteration estimate with an increasing iteration count. \(\square\)

Let \(\tau_3(n,k)\) be the first admissible local stage with
\(q_j(k)\leq1/3\). For every fixed \(0<\epsilon<1/2\), (38) implies, for
all sufficiently large \(b\), \[
 \left\lceil\log_2\frac{V\log3}{1+\epsilon}\right\rceil
 \leq\tau_3(n,k)\leq
 \left\lceil\log_2\frac{V\log3}{1-\epsilon}\right\rceil. \tag{39}
\] For the upper bound, the proposed time lies within
\(\lceil\log_2(6V)\rceil\) because \(\log3/(1-\epsilon)<6\); no
monotonicity assumption about \(h_j(k)\) is used. At the center \(k=n\),
\(V=2n/3\). Thus the earlier fixed-\(j\) central asymptotic is now
controlled uniformly at growing depths, and the local cone-entry time is
\(\log_2n+O(1)\).

\section{Exact certificates and reproducibility}\label{reproducibility-and-verification-audit}


The new verifier uses only the Python standard library and exact integer
arithmetic. Run

\begin{verbatim}
python -B code/verify.py
python -B code/verify.py --baseline
\end{verbatim}

The first command reconstructs and compares every new stored
certificate. The second also reruns the entire earlier verification
suite. \texttt{python -B code/verify.py --write-data} regenerates the new
data rather than comparing it. Assertions are part of the verifier; do
not run Python with \texttt{-O}.

The new files contain all 2047 strict Poisson local certificates, the
540 direct integer-ratio cross-checks, exact structural polynomial
cross-checks, rational constant checks, and recorded execution summaries. For a
separate test of (24)--(29), the verifier constructs every retained
scaled polynomial through depth four from initial indices
\(0\leq m\leq12\), verifies its degree, its leading coefficient against
the independent Poisson iteration, its coefficient-norm bound, and
evaluations at \(n=0,1,5,10,20,50\) against independently convolved
trinomial rows. These finite structural tests support debugging; the
proofs of the formulas for arbitrary indices are the inductions in
Section 7, not an extrapolation from these tests.

The large-cutoff proof uses coefficient-norm and denominator estimates,
not stored coefficient expansions at indices as large as \(B\). The
verifier records the symbolic expressions and checks the modest exponent
inequalities in the derivation; it neither constructs \(\mathcal N\) nor
executes that many exponent checks.

The argument retains the following quantifier and dependency conditions. The support
remains fixed throughout; only artificial computation prefixes are
shortened. The tilted interpolation has an exact sampling identity with
the actual coefficients. Its error is uniform as the tilt approaches
either endpoint. The analytic iteration estimate holds through a
specified growing depth on disks with adequate unit-shift margins. The
half-line seed has strict, finite stopping certificates at every index,
not merely infinitely many nonnegative iterates. The finite boundary
transfer retains all needed neighbors and has a threshold independent of
the subsequent iteration count. Local stopping is converted to global
positivity by a first-failure argument, avoiding circular assumptions.
No real-rootedness or positivity of auxiliary complex variables is
presumed.


The first command is the minimal entry point for the Poisson, polynomial and constant certificates. The optional legacy baseline is included for the finite-exponent comparison and is not a dependency of the proofs of Theorems~A--C. Detailed execution scope is recorded in the accompanying reproducibility document.

The accompanying baseline records \cite{Data} certify the finite exponent set
\[
 \{10,11,\ldots,1000\}\cup\{2000,10000,100000\}.
\]
These finite-exponent results are not presented as new. The cofinite and general-polynomial arguments do not depend on the baseline \(\mathrm{PF}_4\) theorem; they use the initial coefficient formula, ratio dynamics and strict Poisson boundary certificates.

\section{The remaining sharp trinomial threshold}\label{precise-remaining-gap-for-the-requested-target}


Theorem A proves all iterations for every exponent \(n\geq N_0\), and
the accompanying baseline records certify all iterations for their 994 listed
exponents. Thus the still-unresolved portion of (T) is contained in the
finite set \[
 \{n\in\mathbb Z:10\leq n<N_0\}
 \setminus\bigl(\{10,\ldots,1000\}\cup\{2000,10000,100000\}\bigr). \tag{40}
\] This set is finite but astronomically large; it has not been
exhausted computationally. The precise missing theorem is that none of
these exponents is exceptional, or an alternative uniform argument that
supplies all local certificates already for every \(n\geq10\). The
product counterexample of Medina--Straub [1, Example 5.1] prevents treating multiplication by
another copy of \(1+x+x^2\) as a generally valid infinite-log-concavity
induction.

The earlier obstacle was a lack of control when the exponent and
required iteration depth both grow. That obstacle is removed on a common
cofinite tail. What remains is \textbf{the sharp starting threshold},
not a missing all-iteration limit in the new tail theorem. Theorem B
does not settle that sharper question. No complete proof or
counterexample for (T) is claimed.


\paragraph{Attribution and assistance.}
The conjectural context is due to Medina and Straub. The local stopping and shrinking-disk method is prior work, attributed above to Xie, Zhang and Zhang and its cited antecedents. This manuscript was prepared with AI assistance. No global-priority determination, independent human certification or formal proof-assistant verification is claimed.

\begin{thebibliography}{9}
\bibitem{MS} L.~A.~Medina and A.~Straub, \emph{On multiple and infinite log-concavity}, arXiv:1405.1765v1 (7 May 2014); \emph{Annals of Combinatorics} \textbf{20} (2016), 125--138. In particular, Lemma~2.5, Conjecture~5.2 and Examples~5.1 and~5.3. \url{https://arxiv.org/abs/1405.1765v1}.

\bibitem{XZZ} M.~H.~Y.~Xie, C.~X.~T.~Zhang and P.~B.~Zhang, \emph{A unified approach to infinite log-concavity of combinatorial sequences}, arXiv:2610.12237v1 (8 October 2026), Section~2, particularly Lemmas~2.1--2.2, Theorem~2.3 and Corollaries~2.5--2.6. The local closure and shrinking-disk framework is prior work; that source also credits related earlier work of Liu and Mao. The estimates required here are rederived with explicit conservative constants. \url{https://arxiv.org/html/2610.12237v1}.

\bibitem{MCS} P.~R.~W.~McNamara and B.~E.~Sagan, \emph{Infinite log-concavity: developments and conjectures}, arXiv:0808.1065v2; \emph{Advances in Applied Mathematics} \textbf{44} (2010), 1--15. This is background for the factor-cone and baseline determinant framework, rather than an unproved dependency of Theorems~A--C. \url{https://arxiv.org/abs/0808.1065}.

\bibitem{Data} \emph{Supplementary exact computational records for trinomial log-concavity}, 2026. The finite exponent certificates, coefficient records and their standard-library verifier are distributed under \texttt{baseline/}. These computational records are not represented as published literature.
\end{thebibliography}
\end{document}
